\chapter{Voorwoord}

Dit boek heeft een eenvoudige ambitie: één samenhangend natuurkundeboek
dat alles dekt van de kleuterschool tot het einde van de bachelor,
geschreven voor lezers overal ter wereld.

De stijl is beknopt en streng: cursussen opgebouwd uit definities,
voorbeelden, proposities en stellingen, met bewijzen wanneer die op het
gegeven niveau toegankelijk zijn, gevolgd door zorgvuldig gekozen
oefeningen. Resultaten die zonder bewijs worden gegeven, worden
expliciet als toegegeven gemarkeerd. Volledige oplossingen van alle
oefeningen staan achter in het boek.

Elk schooljaar is een deel van het boek, en elk deel is opgesplitst in
hoofdstukken. Hoe vroeger het jaar, hoe jonger de beoogde lezer: meer
illustraties, meer uitgewerkte stappen en zachtere oefeningen.

\section*{Hoe dit boek te lezen}

Uitspraken zijn genummerd per hoofdstuk en kleurgecodeerd: definities in
blauw, stellingen in rood, proposities en verwanten in oranje, en
methodes --- korte recepten voor standaardtaken --- in groen. Oefeningen
zijn genoteerd van ($\star$) voor directe toepassingen tot
($\star\star\star$) voor moeilijkere problemen.

\section*{Bijdragen}

De bron van dit boek staat in een openbare git-repository:
\begin{center}
\url{https://github.com/bvirrion/one-physics-book}
\end{center}
Bijdragen --- nieuwe hoofdstukken, correcties, betere bewijzen, extra
oefeningen --- zijn welkom; zie \texttt{CONTRIBUTING.md} in de root van
de repository.
